En este trabajo estudiamos el teorema de Goodstein, centrándonos en verlo como parte de la aritmética de Peano. También estudiamos a las hidras, objetos matemáticos cuyo comportamiento es también descriptible completamente dentro de la aritmetica de Peano y poseen un teorema analogo al de Goodstein.
Alrededor de estos dos resultados desarrollaremos las teorias de la lógica de primer orden, las funciones computables, la teoría de ordinales, los modelos no estándar de un lenguaje de primer orden, los indicadores de modelos y los conjuntos $\alpha$-grandes para demostrar que ambos resultados son independientes de los axiomas de Peano, es decir, no se pueden demostrar.
Este trabajo es un compilado de resultados cuya intención es dar demostraciones completas y autocontenidas de ambas indemostrabilidades al desarrollar el modelo en donde ambos teoremas no se cumplen, así como explorar las relaciones entre ambos resultados.
El trabajo reúne resultados clásicos de la teoria de modelos, así como una gran variedad de resultados sobre el estudio de la aritmética de Peano, donde brillan los nombres de matemática como Ketonen, Solovay, Kirby y Paris.